% Generated draft: requires mathematical and source review.
\documentclass[11pt]{article}
\usepackage[margin=25mm]{geometry}
\usepackage{fontspec}
\setmainfont{Latin Modern Roman}
\usepackage{amsmath,amssymb,mathtools,unicode-math}
\setmathfont{Latin Modern Math}
\usepackage{enumitem,needspace,xurl,hyperref}
\setlength{\parindent}{0pt}
\setlength{\parskip}{5pt}
\setlength{\emergencystretch}{3em}
\newcommand{\N}{\mathbb{N}}
\newcommand{\Z}{\mathbb{Z}}
\newcommand{\Q}{\mathbb{Q}}
\newcommand{\R}{\mathbb{R}}
\newcommand{\C}{\mathbb{C}}
\newcommand{\F}{\mathbb{F}}
\newcommand{\A}{\mathbb{A}}
\newcommand{\PP}{\mathbb{P}}
\newcommand{\E}{\mathbb{E}}
\newcommand{\eps}{\varepsilon}
\newcommand{\dd}{\,\mathrm{d}}
\begin{document}
\section*{UnsolvedMath ID: 30002002}
\textbf{Source problem code:} catalogue entry
\subsection{Definitions and mathematical statement}
The supplied record does not specify the ground field or the class of Lie algebras denoted by $\mathfrak g$. The following construction gives a precise conditional interpretation without resolving those missing hypotheses.

Suppose that $k$ is a field and that $\mathfrak g$ is a finite-dimensional Lie algebra over $k$. Let
\[
A=\operatorname{Sym}_k(\mathfrak g^*\oplus\mathfrak g^*)
\]
be the coordinate ring of $\mathfrak g\times\mathfrak g$. For each $\lambda\in\mathfrak g^*$, let $f_\lambda\in A$ be the polynomial defined by
\[
f_\lambda(x,y)=\lambda([x,y]).
\]
Define the commutator ideal and commuting scheme by
\[
I_{\mathrm{comm}}=(f_\lambda\mid\lambda\in\mathfrak g^*),
\qquad
\mathcal C_k(\mathfrak g)=\operatorname{Spec}(A/I_{\mathrm{comm}}).
\]
Thus the defining equations are the coordinates of $[x,y]=0$, with their scheme structure retained. A scheme is reduced if its structure sheaf has no nonzero nilpotent sections; for this affine scheme, this means that $A/I_{\mathrm{comm}}$ has no nonzero nilpotent elements.

\textbf{Problem, subject to unresolved hypotheses.} For the ground field $k$ and Lie algebra $\mathfrak g$ intended by the source, determine whether $\mathcal C_k(\mathfrak g)$ is reduced. Equivalently, determine whether
\[
\sqrt{I_{\mathrm{comm}}}=I_{\mathrm{comm}},
\]
or, equivalently, whether for every $a\in A$ and every integer $m\geq 1$,
\[
a^m\in I_{\mathrm{comm}}\ \Longrightarrow\ a\in I_{\mathrm{comm}}.
\]
The record does not determine the characteristic or algebraic closure assumptions on $k$, the structural hypotheses on $\mathfrak g$, or the intended range of quantification over these parameters. In particular, this draft does not assert a conjecture for every finite-dimensional Lie algebra over every field.
\subsection{Short English statement}
The question asks whether the equations saying that two elements of a Lie algebra commute define a scheme with no nonzero nilpotents in its coordinate ring. The field, the hypotheses on the Lie algebra, and the intended scope of the question are missing from the supplied record and require source clarification.
\subsection{Sources}
\begin{itemize}
\item Supplied local problem registry record: permanent UnsolvedMath ID $30002002$, title \emph{Reducedness of the Commuting Scheme}, source problem code \texttt{catalogue entry}, in \emph{Oberwolfach Reports Open Problems}. The record gives the proposed year as 2012 and names no individual proposer.
\item Supplied report citation: \emph{Enveloping Algebras and Geometric Representation Theory} (2012), Oberwolfach Reports, pp.\ 733--809, DOI 10.4171/owr/2012/13. \url{https://doi.org/10.4171/owr/2012/13}. The background describes Oberwolfach Reports as documenting mathematical workshops and their open-problem discussions; no precise problem page or speaker attribution is supplied.
\item Supplied literature source labeled \emph{Oberwolfach Report 13 2012}: \url{https://doi.org/10.4171/owr/2012/13}.
\item Supplied literature source: \emph{On the Normality of Commuting Scheme}, \url{https://arxiv.org/abs/2505.13013}. The supplied assessment says that this 2025 work proves normality only for the ``two-dimensional general-linear case''; its exact hypotheses and the meaning of that phrase are not specified in the record.
\item Literature assessment dated exactly \texttt{2026-08-21}: the supplied curation labels the problem open and describes reducedness of the general commuting scheme as a long-standing conjecture, with the preceding qualification concerning the 2025 work. This is supplied dated curation derived from the OWR material, not an independently verified literature assessment. The record expressly requests independent verification of time-sensitive status and cited claims.
\item The registry labels the statement \texttt{corrected\_verified} and reports a manual semantic, typographic, and LaTeX review of its cleanup. Those supplied labels do not resolve the missing mathematical hypotheses or establish independent mathematical verification of this draft.
\end{itemize}
\end{document}
